\documentclass[11pt, letterpaper]{article}

% --- Packages ---
\usepackage{amsmath, graphicx, xcolor, listings, tcolorbox, booktabs}
\usepackage[left=1in, right=1in, bottom=1in, top=3cm, headheight=2cm, headsep=0.5cm]{geometry}
\usepackage{fancyhdr}

% --- Code Listing Formatting ---
\definecolor{codegreen}{rgb}{0,0.6,0}
\definecolor{codegray}{rgb}{0.5,0.5,0.5}
\definecolor{backcolour}{rgb}{0.95,0.95,0.92}

\lstset{ 
    backgroundcolor=\color{backcolour},   
    commentstyle=\color{codegreen},
    keywordstyle=\color{blue},
    basicstyle=\ttfamily\small,
    breaklines=true,                 
    language=C++
}

% --- Header Configuration ---
\pagestyle{fancy} 
\fancyhf{} 
\renewcommand{\headrulewidth}{0.4pt} 
\lhead{\raisebox{-0.5cm}{\textbf{Universidad de Antioquia}}}
\rhead{\large \textbf{GMAT Targeter Tuning Guide} \\ \normalsize Spaceflight Dynamics 2026-1} 
\cfoot{\thepage} 

\begin{document}

\begin{center}
    {\Large \textbf{Heuristics for Deep-Space Differential Correction}} \\
    \vspace{0.2cm}
    {\large \textit{Parameter Selection for Vary, Propagate, and Achieve}} \\
    \vspace{0.4cm}
    \textit{Prof. Juan | Aerospace Engineering Program}
\end{center}
\hrule
\vspace{0.5cm}

\section*{1. The \texttt{Vary} Command: Controlling the Gradients}
The \texttt{Vary} command calculates the sensitivity matrix (Jacobian) by slightly nudging a variable (\texttt{Perturbation}) and seeing how the outcome changes. It then takes a step toward the solution, limited by \texttt{MaxStep}.

\subsection*{A. True Anomaly (Ignition Timing)}
\begin{lstlisting}
Vary DefaultDC(Voyager.TA = Voyager.TA, {Perturbation = 0.001, Lower = 0, Upper = 360, MaxStep = 5.0});
\end{lstlisting}
\begin{itemize}
    \item \textbf{Perturbation:} $0.001^\circ$ to $0.01^\circ$. Orbit timing is sensitive. If the perturbation is too small, numerical noise ruins the gradient.
    \item \textbf{MaxStep:} $5^\circ$ to $10^\circ$. If you allow the solver to jump $90^\circ$ in a single iteration, the physics become highly non-linear and the solver will diverge. 
\end{itemize}

\subsection*{B. Earth Departure Burn (Trans-Jupiter Injection)}
\begin{lstlisting}
Vary DefaultDC(TJI.Element1 = 7.0, {Perturbation = 1e-04, MaxStep = 0.5});
Vary DefaultDC(TJI.Element2 = 0.0, {Perturbation = 1e-04, MaxStep = 0.1});
Vary DefaultDC(TJI.Element3 = 0.0, {Perturbation = 1e-04, MaxStep = 0.1});
\end{lstlisting}
\begin{itemize}
    \item \textbf{Element 1 (Velocity):} Give it a strong initial guess (e.g., $7.0$ km/s for Jupiter). Set \texttt{MaxStep = 0.5} km/s. Do not let it jump by $5$ km/s at a time.
    \item \textbf{Element 2 \& 3 (Normal/Binormal):} Initial guess $0.0$. Keep \texttt{MaxStep = 0.1} km/s. Out-of-plane maneuvers are incredibly expensive; keep the solver from wasting fuel here.
\end{itemize}

\subsection*{C. Deep Space TCM (Trajectory Correction Maneuver)}
\begin{lstlisting}
Vary DefaultDC(TCM.Element1 = 0.0, {Perturbation = 1e-05, MaxStep = 0.05});
\end{lstlisting}
\begin{itemize}
    \item \textbf{Perturbation:} $10^{-5}$ km/s. Deep space burns are micro-adjustments.
    \item \textbf{MaxStep:} $0.01$ to $0.05$ km/s. A $1$ m/s change millions of kilometers away radically alters the Jupiter arrival point.
\end{itemize}

% ---------------------------------------------------------
\section*{2. The \texttt{Achieve} Command: Defining Success}
A common student mistake is demanding absolute mathematical perfection (e.g., $0.0000001$ tolerance). In an $N$-body universe, perfection is impossible and forces the solver into an infinite loop.

\subsection*{A. Earth Escape Asymptote}
\begin{lstlisting}
Achieve DefaultDC(Voyager.EarthMJ2000Eq.C3Energy = 85.5, {Tolerance = 0.05});
Achieve DefaultDC(Voyager.EarthMJ2000Eq.DLA = 15.2, {Tolerance = 0.05});
Achieve DefaultDC(Voyager.EarthMJ2000Eq.RLA = 120.4, {Tolerance = 0.05});
\end{lstlisting}
\begin{itemize}
    \item \textbf{Tolerances:} $0.01$ to $0.1$. The Earth departure burn does not need to be microscopic in its accuracy because the deep-space TCM will clean up the residual errors. Allow the solver to say "close enough."
\end{itemize}

\subsection*{B. Jupiter B-Plane Targeting}
\begin{lstlisting}
Achieve DefaultDC(Voyager.Jupiter.BdotT = 400000, {Tolerance = 50.0});
Achieve DefaultDC(Voyager.Jupiter.BdotR = 150000, {Tolerance = 50.0});
\end{lstlisting}
\begin{itemize}
    \item \textbf{Tolerances:} $10$ km to $50$ km. You are throwing a dart from Earth and trying to hit a moving target at Jupiter 3 years later. A $50$ km tolerance on a target that is $800$ million kilometers away is phenomenal engineering. 
\end{itemize}

% ---------------------------------------------------------
\section*{3. The \texttt{Propagate} Command: The Fail-Safes}
Never propagate to a spatial condition (\texttt{Periapsis} or \texttt{Apoapsis}) without including a temporal fail-safe (\texttt{ElapsedDays}).

\begin{lstlisting}
% Earth Escape (Clearing the Sphere of Influence)
Propagate SunProp(Voyager) {Voyager.Earth.Distance = 1000000, Voyager.ElapsedDays = 30};

% Jupiter Transit
Propagate SunProp(Voyager) {Voyager.Jupiter.Periapsis, Voyager.ElapsedDays = 1500};
\end{lstlisting}
\begin{itemize}
    \item \textbf{Jupiter Transit:} Earth to Jupiter takes roughly $1,000$ days. Setting the fail-safe to $1,500$ gives the targeter room to explore "slow" trajectories, but forcibly stops the simulation before the spacecraft flies off the end of the JPL Ephemeris calendar (preventing the \texttt{leDE1941.405} crash).
\end{itemize}

\end{document}